% Suggested insertion following the discussion of Popper's C and C^+ predictors.
% Requires csquotes and biblatex (or adapt \autocite commands to your citation package).

The significance of Popper's argument is therefore stronger than the observation that a sufficiently complicated machine may fail to predict itself. The obstruction is logical rather than computational. If the information-bearing state of the predictor $C$ is itself among the facts to be represented, then supplying $C$ with a complete description of that state changes the very state described. The description is therefore not stable under its own incorporation into the system. From the standpoint of $C^{+}$ there need be no corresponding indeterminacy: the evolution of the combined system $C+E$ may be described deterministically, including the correlations established by their interaction. What fails for $C$ is the possibility of placing its own state, its prediction, and the subsequent $C$--$E$ interaction within a single, simultaneously valid description available to $C$ itself. Popper's argument should not, by itself, be taken to prove the quantum formalism; rather, it identifies a self-referential obstruction to a global classical description. This is strikingly close to the modern characterization of contextuality as an obstruction to assembling locally consistent descriptions into a single global assignment, a connection which has also been related explicitly to logical paradox and cyclic self-reference \autocite{Abramsky2015Contextuality,Dourdent2020Godelian}.

This observation provides a route from Popper's argument to negative probability. Ordinary Kolmogorov probability assumes a single sample space on which all relevant events can be represented by a non-negative joint probability measure. Fine's theorem shows, in the Bell--CHSH setting, how substantive this assumption is: the existence of an appropriate non-negative joint distribution is equivalent to satisfaction of the Bell inequalities \autocite{Fine1982Hidden}. Conversely, where the observable marginals cannot be embedded in such a distribution, one may retain their probabilistic bookkeeping only by abandoning some classical requirement. One particularly conservative extension is to retain normalization and the required marginals while relaxing positivity, thereby replacing probability by a signed or quasi-probability distribution. Halliwell and Yearsley make precisely this connection between negative probabilities and the failure of a viable positive joint distribution \autocite{HalliwellYearsley2013}. More generally, Spekkens showed, for an operationally defined class of representations, a close equivalence between the possibility of an everywhere non-negative quasiprobability representation and the existence of a noncontextual ontological model \autocite{Spekkens2008Negativity}. Negativity is therefore not merely an exotic numerical device: it can be understood as the mathematical residue left when locally meaningful probabilistic descriptions cannot be combined into a single noncontextual positive one.

Popper's predictor thus suggests an epistemological origin for precisely the kind of obstruction for which quasi-probability provides a natural representation. Self-reference does not logically entail that an individual observable event possesses a negative probability; observable probabilities remain non-negative. Rather, it prevents the predictor from consistently treating all of the counterfactual and self-descriptive propositions required for its prediction as events in one stable positive probability space. If we nevertheless insist on representing these incompatible partial descriptions within a common calculus, signed probabilities provide a minimal extension of classical probability. This still does not by itself explain why quantum mechanics requires \emph{complex} probability amplitudes. Negativity supplies the possibility of cancellation, but complex amplitudes additionally carry phase and hence the structure required for interference. The proposed chain of explanation is therefore not that Popper's theorem alone derives quantum mechanics, but that self-reference motivates the failure of globally positive classical probability; contextuality makes that failure precise; signed quasi-probability represents the resulting obstruction; and the additional empirical requirement of interference motivates the further passage from real signed weights to complex amplitudes.
